\section{Adjoints in moment coordinates}\label{sec:identities}

We retain the probability measures
\[
 d\mu(x)=Z^{-1}e^{-V(x)}\,dx,
 \qquad d\nu(z)=e^{-\varphi(z)}\,dz,
 \qquad x=\nabla\varphi(z),
\]
and the smooth diffeomorphism constructed above.  Put
\[
 \tau=D_z^2\varphi,\qquad M=\tau^{-1}.
\]
In target coordinates, these symbols denote the same matrix fields
evaluated at $z=(\nabla\varphi)^{-1}(x)$.
Expectation under either measure, with functions identified by this map,
is denoted by $\E$.
Throughout the coordinate calculations, repeated indices are summed
from $1$ to $n$, and $Q:L=\tr(Q^TL)$ denotes matrix contraction.

We first construct the compact test functions used in the variance
argument and establish an identity that will also identify the
obstruction to their density.  Write
\[
 \partial_k=\partial_{z_k},\qquad
 \partial_k^*=x_k-\partial_k.
\]
The latter is the adjoint of $\partial_k$ in $L^2(\nu)$ because
$\partial_k\varphi=x_k$.  In target coordinates, write $\delta_\mu$
for negative weighted divergence:
\[
 \delta_\mu W
 =\sum_j\bigl(V_{x_j}W_j-\partial_{x_j}W_j\bigr).
\]
On a matrix, $\delta_\mu$ acts on each row, producing a vector.
For a smooth row field $P_{ak}$ and a compact test $g$, the chain rule
and adjunction give
\[
 \E g\,\partial_k^*P_{ak}
 =\E(\partial_k g)P_{ak}
 =\E g_{x_i}\tau_{ik}P_{ak}.
\]
Consequently
\begin{equation}\label{eq:row-adjoint}
 \partial_k^*P_{ak}=(\delta_\mu(P\tau))_a.
\end{equation}
Taking $P$ to be the identity matrix also gives $\delta_\mu\tau=x$.
This is Fathi's moment-map construction of a Stein kernel
\cite[Theorem~2.3 of the cited arXiv version]{Fathi2019}: the matrix
field $\tau$ expresses integration by parts against $x$ through
$\E x_j g=\E\tau_{jk}g_{x_k}$.

Thus the differential operator
\begin{equation}\label{eq:operator-A}
 Af=\delta_\mu(\tau\nabla_x f)
   =\partial_k^*f_{x_k}
   =x\cdot\nabla_x f-\tau:D_x^2f
\end{equation}
is formally symmetric, with
\begin{equation}\label{eq:A-form}
 \E gAf=\E\nabla_xg\cdot\tau\nabla_xf
\end{equation}
for compact smooth tests.  These formulas involve no operator-domain
assertion. This is the negative of the moment-Hessian diffusion
generator studied by Klartag \cite[Section~6]{KlartagArxiv2013};
for our sign convention, $Ax_j=x_j$.

For $f\in C_c^\infty(\R_x^n)$, define
\begin{equation}\label{eq:test-fields}
 h=D_x^2f,\qquad P=\tau h,\qquad B=\tau h\tau,
 \qquad W_a=\partial_k^*P_{ak},\qquad u=\delta_\mu W.
\end{equation}
All these fields have compact support in both coordinates: the inverse
image of a compact target set is its compact image under the continuous
inverse diffeomorphism.  By \eqref{eq:row-adjoint},
$W=\delta_\mu B$.  For every smooth function $F$ in target coordinates,
two adjunctions give
\begin{equation}\label{eq:test-pairing}
 \E uF=\E W\cdot\nabla_xF=\E B:D_x^2F.
\end{equation}
The fields are compactly supported, so $F$ need not be.
Thus the same pairing can be estimated through either the first or the
second derivatives of $F$.
To identify the range of these tests, we also use the gradient identity
\begin{equation}\label{eq:W-gradient}
 W=\nabla_z\bigl[((A-1)f)(x(z))\bigr]
   =\tau\nabla_x(A-1)f,
 \qquad u=A(A-1)f.
\end{equation}
To check it directly, put
$C_{ijk}=\partial_{z_i}\partial_{z_j}\partial_{z_k}\varphi$.
This tensor is fully symmetric, and expansion gives
\[
 W_a=x_k\tau_{ai}f_{x_ix_k}
       -C_{aik}f_{x_ix_k}
       -\tau_{ai}\tau_{jk}f_{x_ix_kx_j}.
\]
Differentiating $(A-1)f$ in $z_a$ gives exactly these terms: the
first-gradient terms from $x\cdot\nabla_xf$ and $-f$ cancel.
The formula for $u$ follows from \eqref{eq:operator-A}.
In particular, $D_zW$ is symmetric.

The following identity is the moment-coordinate specialization of
the transport-Hessian Bochner formula of Kolesnikov, Livshyts and
Rotem \cite[Proposition~4.8]{KLR2026}. We give its direct derivation
in the present notation. Two compact integrations by parts give
\begin{equation}\label{eq:hessian}
 \begin{aligned}
 \E(Af)(A-1)f
 &=\E\nabla_x f\cdot\tau\nabla_x(A-1)f
   =\E\nabla_x f\cdot W\\
 &=\E h:B
   =\E\bigl\|\tau^{1/2}D_x^2f\,\tau^{1/2}\bigr\|_{\mathrm{HS}}^2.
 \end{aligned}
\end{equation}
Here $\|Q\|_{\mathrm{HS}}^2=\tr(Q^TQ)$.
The differential expression $A$ satisfies $A1=0$ and $Ax_j=x_j$,
so $A(A-1)$ annihilates every affine function.
Section~\ref{sec:density} uses \eqref{eq:hessian} to prove that these
are the only $L^2$ obstructions to density. Centering removes constants,
and evenness removes linear functions.
